% ===================================================================
%  ਸਾਰਣੀ ਨੰ. 3 — ਬਚਪਨ ਤੇ ਜਵਾਨੀ।
%  Traduction 2026 d'apres texto/io, controlee sur texto/fr, texto/en
%  et texto/hi. Consigne : voir texto/pa/10-sarni-01.tex.
%
%  Deux scenes, et l'ido subdivise beaucoup plus que le francais :
%  huit des intertitres n'ont pas de vis-a-vis francais. Le gourmoukhi
%  suit l'ido, comme les dix autres traductions.
% ===================================================================

%%K t03-apar-1 apar
\VUsaut{3.0mm}
\VUtitre{15.0pt}{60}{ਸਾਰਣੀ ਨੰ. 3}
\VUsaut{1.6mm}
\VUtitre{11.4pt}{0}{ਬਚਪਨ ਤੇ ਜਵਾਨੀ।}
\VUsaut{3.4mm}

%%K t03-apar-2 apar
(ਸਾਰਣੀਆਂ 3 ਤੇ 4 ਇਸ ਤਰ੍ਹਾਂ ਰਚੀਆਂ ਗਈਆਂ ਹਨ ਕਿ ਚਾਰ \VUgras{ਪਰਿਵਾਰਕ ਦ੍ਰਿਸ਼ਾਂ} ਵਿਚ
\VUgras{ਰੁੱਤਾਂ}, \VUgras{ਉਮਰਾਂ}, \VUgras{ਪਰਿਵਾਰ}, \VUgras{ਕੱਪੜੇ} ਤੇ
\VUgras{ਸਿਹਤ} ਦੀ ਮੁਢਲੀ ਸ਼ਬਦਾਵਲੀ ਸਾਹਮਣੇ ਆ ਜਾਵੇ।)
\VUsaut{4.0mm}

%%K t03-tit-1 sub
\VUcentre{12.0pt}{0}{\textit{ਪਹਿਲਾ ਦ੍ਰਿਸ਼।}}
\VUsaut{3.0mm}

%%K t03-tit-2 sub
\VUpk{10.2pt}{40}{ਪਿੰਡ ਵਿਚ ਇਕ ਬਪਤਿਸਮਾ।}
\VUsaut{2.4mm}

%%K t03-tit-3 sub
\VUpk{10.2pt}{40}{ਬਹਾਰ।}
\VUsaut{3.0mm}

%%K t03-c1-01-1 p
1. --- ਅਸੀਂ \VUgras{ਬਹਾਰ} ਵਿਚ ਹਾਂ, \VUgras{ਮਾਰਚ}, \VUgras{ਅਪ੍ਰੈਲ} ਜਾਂ
\VUgras{ਮਈ} ਦੇ ਮਹੀਨੇ ਵਿਚ, \VUgras{ਹਫ਼ਤੇ} ਦੇ ਕਿਸੇ ਦਿਨ --- \VUgras{ਸੋਮਵਾਰ},
\VUgras{ਮੰਗਲਵਾਰ} ਜਾਂ \VUgras{ਬੁੱਧਵਾਰ}। ਸਵੇਰ ਦੇ \VUgras{ਦਸ ਵਜੇ} ਹਨ।

%%K t03-c1-02-1 p
2. --- \VUgras{ਮੌਸਮ} ਸੋਹਣਾ ਹੈ, \VUgras{ਹਵਾ} ਨਿਖਰੀ ਹੋਈ। ਸੂਰਜ ਚਮਕ ਰਿਹਾ ਹੈ। ਕੁਝ
ਨਿੱਕੇ ਨਿੱਕੇ \VUgras{ਬੱਦਲ} \textsuperscript{(2)} ਨੀਲੇ \VUgras{ਅਸਮਾਨ}
\textsuperscript{(1)} ਵਿਚ ਉੱਠ ਰਹੇ ਹਨ।

%%K t03-c1-03-1 p
3. --- \VUgras{ਰੁੱਖਾਂ} ਉੱਤੇ ਹਾਲੇ \VUgras{ਪੱਤੇ} ਤਕਰੀਬਨ ਨਹੀਂ। ਪਤਲੇ
\VUgras{ਸਫ਼ੈਦੇ} \textsuperscript{(3)} ਤੇ ਉੱਚੇ \VUgras{ਚਨਾਰ} \textsuperscript{(17)}
ਉੱਤੇ ਹੁਣ ਤਕ ਕਲੀਆਂ ਹੀ ਹਨ।

%%K t03-c1-04-1 p
4. --- \VUgras{ਅਬਾਬੀਲਾਂ} \textsuperscript{(4)} ਮੁੜ ਆਈਆਂ ਹਨ, ਤੇ ਖ਼ੁਸ਼ੀ ਨਾਲ
ਕੂਕਦੀਆਂ ਹੋਈਆਂ ਉਸ \VUgras{ਚੋਟੀ} \textsuperscript{(7)} ਦੇ ਚੁਫੇਰੇ ਚੱਕਰ ਕੱਢ ਰਹੀਆਂ
ਹਨ ਜੋ \VUgras{ਘੰਟਾ-ਘਰ} \textsuperscript{(5)} ਦੀ ਹੈ, ਤੇ ਉਹ ਘੰਟਾ-ਘਰ ਪਿੰਡ ਦੇ
\VUgras{ਗਿਰਜੇ} \textsuperscript{(6)} ਉੱਤੇ ਖੜਾ ਹੈ।

%%K t03-tit-4 sub
\VUsaut{3.4mm}
\VUpk{10.2pt}{40}{ਗਿਰਜਾ।}
\VUsaut{3.0mm}

%%K t03-c1-05-1 p
5. --- ਇਹ ਗਿਰਜਾ ਬਹੁਤ ਸਾਦਾ ਹੈ, ਆਪਣੇ ਵੱਡੇ \VUgras{ਦਰਵਾਜ਼ੇ} \textsuperscript{(13)}
ਤੇ ਆਪਣੀ ਵੱਡੀ \VUgras{ਘੜੀ} \textsuperscript{(8)} ਨਾਲ। \VUgras{ਨਿੱਕੀ ਸੂਈ}
\textsuperscript{(9)} X ਦੇ ਹਿੰਦਸੇ ਉੱਤੇ ਹੈ: ਉਹ \VUgras{ਘੰਟੇ} ਦੱਸਦੀ ਹੈ।
\VUgras{ਵੱਡੀ ਵਾਲੀ} \textsuperscript{(10)} XII (ਦੁਪਹਿਰ) ਉੱਤੇ ਹੈ; ਉਹ
\VUgras{ਮਿੰਟ} ਦੱਸਦੀ ਹੈ।

%%K t03-c1-06-1 p
6. --- ਮੁਨਾਰੇ ਵਿਚ \VUgras{ਘੰਟੇ} \textsuperscript{(11)} ਖ਼ੁਸ਼ੀ ਨਾਲ ਵੱਜ ਰਹੇ ਹਨ। ਛੱਤ
ਦੇ ਸਿਰੇ ਉੱਤੇ ਪੱਥਰ ਦੀ ਇਕ ਨਿੱਕੀ \VUgras{ਸਲੀਬ} \textsuperscript{(12)} ਖੜੀ ਹੈ।

%%K t03-c1-07-1 p
7. --- ਗਿਰਜੇ ਵਿਚ ਸਿਰਫ਼ ਵੱਡਾ \VUgras{ਦਰਵਾਜ਼ਾ} \textsuperscript{(13)} ਹੀ ਨਹੀਂ, ਬਗ਼ਲ
ਵਿਚ ਇਕ ਨਿੱਕਾ ਬੂਹਾ ਵੀ ਹੈ। ਉਸੇ ਬੂਹੇ ਵਿਚੋਂ \VUgras{ਪਾਦਰੀ} \textsuperscript{(15)},
ਆਪਣੇ \VUgras{ਚੋਗ਼ੇ} ਵਿਚ, \VUgras{ਕੱਪੜਿਆਂ ਵਾਲੇ ਕਮਰੇ} \textsuperscript{(14)} ਤੋਂ
ਨਿਕਲ ਕੇ \VUgras{ਪਾਦਰੀ ਦੇ ਘਰ} \textsuperscript{(19)} ਵੱਲ ਜਾਂਦੇ ਹਨ।

%%K t03-c1-08-1 p
8. --- ਉਨ੍ਹਾਂ ਕੋਲ ਇਹ ਰਹੀ \VUgras{ਦੁੱਧ ਵਾਲੀ} \textsuperscript{(24)} ਆਪਣੇ
\VUgras{ਖੋਤੇ} \textsuperscript{(23)} ਨਾਲ। ਉਹ ਸ਼ਹਿਰੋਂ ਮੁੜ ਰਹੀ ਹੈ, ਜਿੱਥੇ ਬੱਚਿਆਂ ਲਈ
ਆਪਣਾ ਚੰਗਾ ਦੁੱਧ ਪਹੁੰਚਾ ਆਈ ਹੈ।

%%K t03-tit-5 sub
\VUsaut{4.0mm}
\VUpk{10.2pt}{40}{ਫਲਾਂ ਦਾ ਬਾਗ਼।}
\VUsaut{3.4mm}

%%K t03-c1-09-1 p
9. --- ਉਸਤਾਦ ਅਲੀਬੋਰੋਂ (ਖੋਤਾ) ਇਸ ਸਵੇਰ ਦੀ ਦੌੜ ਤੋਂ ਬਹੁਤ ਖ਼ੁਸ਼ ਹੈ। ਉਹ ਗੁਆਂਢ ਦੇ
\VUgras{ਬਾਗ਼} ਦੇ ਉਨ੍ਹਾਂ \VUgras{ਫੁੱਲਾਂ} \textsuperscript{(20)} ਤੋਂ ਮਹਿਕੀ ਹਵਾ
ਆਨੰਦ ਨਾਲ ਅੰਦਰ ਖਿੱਚਦਾ ਹੈ ਜੋ \VUgras{ਫਲਾਂ ਵਾਲੇ ਰੁੱਖਾਂ} \textsuperscript{(18)} ਨੂੰ
ਢਕੀ ਬੈਠੇ ਹਨ। ਉਹ \VUgras{ਆੜੂ ਦੇ ਰੁੱਖ} \textsuperscript{(18)} ਤੇ
\VUgras{ਖ਼ੁਬਾਨੀ ਦੇ ਰੁੱਖ} \textsuperscript{(19)} ਬਿਲਕੁਲ ਗੁਲਾਬੀ ਤੇ ਚਿੱਟੇ ਹਨ, ਤੇ
ਚੰਗੀ ਫ਼ਸਲ ਦੀ ਆਸ ਬੰਨ੍ਹਾਉਂਦੇ ਹਨ।

%%K t03-c1-10-1 p
10. --- ਇਸ ਦੌਰਾਨ ਨਿੱਕੀਆਂ \VUgras{ਮੱਖੀਆਂ} \textsuperscript{(22)}, ਜੋ
\VUgras{ਛੱਤੇ} \textsuperscript{(21)} ਤੋਂ ਆਉਂਦੀਆਂ ਹਨ, ਉੱਥੇ ਜਾ ਕੇ ਗੁਣਗੁਣਾਉਂਦੀਆਂ
ਆਪਣਾ ਮਿੱਠਾ \VUgras{ਸ਼ਹਿਦ} ਇਕੱਠਾ ਕਰਦੀਆਂ ਹਨ।

%%K t03-c1-11-1 p
11. --- ਪਰ ਇਹ ਕੀ ਹੋ ਰਿਹਾ ਹੈ? ਪਿੰਡ ਦੇ ਬੱਚੇ ਦੌੜਦੇ ਹੋਏ ਆ ਰਹੇ ਹਨ ਤੇ ਡਰੀਆਂ ਹੋਈਆਂ
\VUgras{ਬੱਤਖਾਂ} \textsuperscript{(25)} ਨੂੰ ਖਿਲਾਰ ਰਹੇ ਹਨ। ਇਹ ਉਹ ਜਲੂਸ ਹੈ ਜੋ ਮੁਨਸ਼ੀ
ਦੇ ਪੁੱਤਰ ਦੇ \VUgras{ਬਪਤਿਸਮੇ} ਤੋਂ ਬਾਅਦ ਗਿਰਜੇ ਵਿਚੋਂ ਨਿਕਲ ਰਿਹਾ ਹੈ।

%%K t03-c1-12-1 p
12. --- ਨਿੱਕਾ ਜੇਮਜ਼, ਆਪਣੇ \VUgras{ਪੰਘੂੜੇ} \textsuperscript{(26)} ਵਿਚ, ਘੰਟੀਆਂ ਦੀ
ਖ਼ੁਸ਼ ਆਵਾਜ਼ ਨਾਲ ਜਾਗ ਪੈਂਦਾ ਹੈ, ਤੇ ਉਸ ਦੀ ਭੈਣ ਮੈਡਲਿਨ, ਜੋ ਜਲੂਸ ਵੇਖਣਾ ਚਾਹੁੰਦੀ ਹੈ,
ਆਪਣੀ ਮਾਂ ਨੂੰ ਆਖਦੀ ਹੈ ਕਿ ਆ ਕੇ ਦਹਿਲੀਜ਼ ਉੱਤੇ ਹੀ ਉਸ ਦੇ ਵਾਲ ਸੰਵਾਰ ਦੇਵੇ ਤੇ ਉਸ ਨੂੰ
ਕੱਪੜੇ ਪੁਆ ਦੇਵੇ।

%%K t03-c1-13-1 p
13. --- ਉਸ ਦੀ ਮਾਂ ਮੰਨ ਜਾਂਦੀ ਹੈ। ਉਹ ਆਪਣਾ \VUgras{ਸਿਰ ਦਾ ਕੱਪੜਾ}
\textsuperscript{(28)}, ਆਪਣੀ \VUgras{ਚੋਲੀ} \textsuperscript{(29)} ਤੇ ਆਪਣਾ
\VUgras{ਪਰਨਾ} \textsuperscript{(30)} ਪਾ ਕੇ ਬੂਹੇ ਦੇ ਸਾਹਮਣੇ ਇਕ ਨਿੱਕੀ ਕੁਰਸੀ ਉੱਤੇ ਆ
ਬਹਿੰਦੀ ਹੈ। ਮੈਡਲਿਨ ਉੱਤੇ ਸਿਰਫ਼ ਉਸ ਦਾ \VUgras{ਪੇਟੀਕੋਟ} \textsuperscript{(31)} ਤੇ ਉਸ
ਦੀ \VUgras{ਕੁੜਤੀ} \textsuperscript{(32)} ਹੈ।

%%K t03-c1-14-1 p
14. --- ਉਹ ਕਿੰਨੀ ਖ਼ੁਸ਼ ਹੈ! ਉਹ ਆਪਣੇ ਪਿਆਰੇ ਫੁੱਲ ਭੁੱਲ ਜਾਂਦੀ ਹੈ — ਆਪਣੇ
\VUgras{ਲੌਂਗ ਦੇ ਫੁੱਲ} \textsuperscript{(34)}, ਜਿਨ੍ਹਾਂ ਉੱਤੇ \VUgras{ਤਿੱਤਲੀਆਂ}
\textsuperscript{(35)} ਬਹਿੰਦੀਆਂ ਹਨ, ਆਪਣੇ \VUgras{ਟਿਊਲਿਪ} \textsuperscript{(36)},
ਆਪਣੀਆਂ \VUgras{ਵਾਦੀ ਦੀਆਂ ਸੋਸਨਾਂ} \textsuperscript{(37)} ਜੋ \VUgras{ਗਮਲਿਆਂ}
\textsuperscript{(38)} ਵਿਚ ਹਨ, ਤੇ ਉਹ ਗਮਲੇ \VUgras{ਕੰਧ} \textsuperscript{(33)} ਉੱਤੇ
ਰੱਖੇ ਹਨ, ਤੇ ਆਪਣੇ \VUgras{ਗੁਲਾਬ} \textsuperscript{(39)} ਜੋ ਇੰਨੀ ਸੋਹਣੀ ਵਾੜ
ਬਣਾਉਂਦੇ ਹਨ। ਉਹ ਜਲੂਸ ਦੀਆਂ ਬੀਬੀਆਂ ਦੇ ਕੀਮਤੀ ਕੱਪੜੇ ਤੱਕ ਰਹੀ ਹੈ।

%%K t03-c1-15-1 p
15. --- ਪਿੰਡ ਦੇ ਮੁੰਡੇ ਕੱਪੜਿਆਂ ਵੱਲ ਘੱਟ ਧਿਆਨ ਦਿੰਦੇ ਹਨ। ਉਹ \VUgras{ਮਿਠਾਈਆਂ}
\textsuperscript{(55)} ਜਾਂ \VUgras{ਸਿੱਕੇ} \textsuperscript{(52)} ਲੈਣ ਲਈ ਅੱਗੇ
ਲਪਕਦੇ ਤੇ ਧੱਕਮ-ਧੱਕਾ ਕਰਦੇ ਹਨ। ਉਨ੍ਹਾਂ ਵਿਚੋਂ ਇਕ ਰੋ \textsuperscript{(40)} ਰਿਹਾ ਹੈ।

%%K t03-c1-16-1 p
16. --- ਦੂਜੇ ਨੇ \VUgras{ਕੁੱਤੇ} \textsuperscript{(42)} ਦਾ ਪੰਜਾ ਮਿੱਧ ਦਿੱਤਾ ਹੈ, ਜੋ
ਕੂੰ-ਕੂੰ ਕਰਦਾ ਨੱਸ ਜਾਂਦਾ ਹੈ ਤੇ \VUgras{ਕੁਕੜੀ} \textsuperscript{(43)} ਤੇ ਉਸ ਦੇ
\VUgras{ਚੂਚਿਆਂ} \textsuperscript{(44)} ਨੂੰ ਉਡਾ ਦਿੰਦਾ ਹੈ।

%%K t03-tit-6 sub
\VUsaut{5.0mm}
\VUpk{10.2pt}{40}{ਬਪਤਿਸਮੇ ਦਾ ਜਲੂਸ।}
\VUsaut{4.0mm}

%%K t03-c1-17-1 p
17. --- ਜਲੂਸ ਦੇ ਅੱਗੇ ਅੱਗੇ \VUgras{ਦਾਈ} \textsuperscript{(45)} ਤੁਰ ਰਹੀ ਹੈ। ਉਸ ਦੇ
ਸਿਰ ਉੱਤੇ ਲੰਮੀਆਂ ਫੀਤੀਆਂ ਵਾਲੀ ਸੋਹਣੀ \VUgras{ਟੋਪੀ} \textsuperscript{(46)} ਹੈ। ਉਹ
\VUgras{ਨਵੇਂ ਜੰਮੇ ਬਾਲ} \textsuperscript{(47)} ਨੂੰ ਚੁੱਕੀ ਹੋਈ ਹੈ, ਜੋ ਸਾਰੇ ਦਾ ਸਾਰਾ
\VUgras{ਲੇਸ} \textsuperscript{(48)} ਵਿਚ ਲਪੇਟਿਆ ਹੋਇਆ ਹੈ।

%%K t03-c1-18-1 p
18. --- ਦਾਈ ਦੇ ਪਿੱਛੇ \VUgras{ਧਰਮ-ਪਿਤਾ} \textsuperscript{(49)} ਤੇ
\VUgras{ਧਰਮ-ਮਾਤਾ} \textsuperscript{(53)} ਆਉਂਦੇ ਹਨ। ਉਹ ਮਿਠਾਈਆਂ ਤੇ ਸਿੱਕੇ ਵੰਡ ਰਹੇ
ਹਨ। ਧਰਮ-ਪਿਤਾ ਦੇ ਹੱਥ ਵਿਚ \VUgras{ਦਸਤਾਨੇ} \textsuperscript{(51)} ਤੇ ਸਿਰ ਉੱਤੇ
\VUgras{ਉੱਚੀ ਹੈਟ} \textsuperscript{(50)} ਹੈ। ਧਰਮ-ਮਾਤਾ ਹੱਥ ਵਿਚ ਇਕ ਸੋਹਣਾ
\VUgras{ਫੁੱਲਾਂ ਦਾ ਗੁਲਦਸਤਾ} \textsuperscript{(54)} ਚੁੱਕੀ ਹੋਈ ਹੈ।

%%K t03-c1-19-1 p
19. --- ਉਨ੍ਹਾਂ ਦੇ ਪਿੱਛੇ ਸਾਨੂੰ ਬਾਲ ਦੇ \VUgras{ਚਾਚਾ} \textsuperscript{(56)} ਤੇ
\VUgras{ਚਾਚੀ} \textsuperscript{(58)} ਵਿਖਾਈ ਦਿੰਦੇ ਹਨ। ਚਾਚੀ ਦੇ ਸਿਰ ਉੱਤੇ ਸਜੀਲੀ
\VUgras{ਹੈਟ} \textsuperscript{(59)} ਹੈ, ਚਾਚਾ ਉੱਤੇ ਕਾਲਾ \VUgras{ਲੰਮਾ ਕੋਟ}
\textsuperscript{(57)} ਤੇ ਚਿੱਟੀ ਵਾਸਕਟ। ਫਿਰ \VUgras{ਚਚੇਰਾ ਭਰਾ}
\textsuperscript{(60)} ਤੇ ਉਸ ਦੀ \VUgras{ਚਚੇਰੀ ਭੈਣ} \textsuperscript{(61)} ਆਉਂਦੇ ਹਨ।
ਉਹ ਨਵੇਂ ਜੰਮੇ ਬਾਲ ਦੀ \VUgras{ਭੈਣ} \textsuperscript{(62)}, ਨਿੱਕੀ ਬਰਥਾ, ਦਾ ਹੱਥ ਫੜੀ
ਹੋਈ ਹੈ।

%%K t03-c1-20-1 p
20. --- ਬਰਥਾ ਦਾ ਦੂਜਾ \VUgras{ਭਰਾ} \textsuperscript{(63)} ਪਿੱਛੇ ਤੁਰ ਰਿਹਾ ਹੈ। ਉਹ
ਇਕ \VUgras{ਰਿਸ਼ਤੇਦਾਰ} \textsuperscript{(64)} ਨੂੰ ਹੱਥ ਦਿੱਤਾ ਹੋਇਆ ਹੈ। ਪਰਿਵਾਰ ਦੇ
\VUgras{ਮਿੱਤਰ} \textsuperscript{(65)} ਉਨ੍ਹਾਂ ਤੋਂ ਬਾਅਦ ਆਉਂਦੇ ਹਨ।

%%K t03-tit-7 sub
\VUsaut{4.6mm}
\VUpk{10.2pt}{40}{ਤਮਾਸ਼ਬੀਨ।}
\VUsaut{3.4mm}

%%K t03-c1-21-1 p
21. --- ਜਲੂਸ ਕੋਲ ਇਹ ਰਿਹਾ ਇਕ \VUgras{ਮੰਗਤਾ} \textsuperscript{(66)}, ਆਪਣੀ
\VUgras{ਬੈਸਾਖੀ} \textsuperscript{(67)} ਦੇ ਸਹਾਰੇ। ਉਹ ਖ਼ੈਰਾਤ ਮੰਗ ਰਿਹਾ ਹੈ।

%%K t03-c1-22-1 p
22. --- ਦੂਜੇ ਪਾਸੇ ਇਕ ਬਹੁਤ \VUgras{ਸ਼ਿਸ਼ਟ} \textsuperscript{(68)} ਮੁੰਡਾ ਹੈ। ਉਹ ਆਪਣੀ
ਟੋਪੀ ਚੁੱਕ ਕੇ ਜਾਣ-ਪਛਾਣ ਵਾਲਿਆਂ ਨੂੰ \VUgras{« ਨਮਸਕਾਰ »} ਆਖਦਾ ਹੈ।

%%K t03-c1-23-1 p
23. --- ਪਿੰਡ ਵਾਲੇ ਬਪਤਿਸਮੇ ਦਾ ਜਲੂਸ ਨਿਕਲਦਾ ਵੇਖਣ ਆਪਣੇ ਘਰਾਂ ਤੋਂ ਬਾਹਰ ਆ ਗਏ ਹਨ।
ਸਾਨੂੰ ਆਪਣੀ \VUgras{ਸੂਤੀ ਟੋਪੀ} ਪਾਈ ਇਕ \VUgras{ਪੇਂਡੂ ਬੰਦਾ} \textsuperscript{(69)}
ਵਿਖਾਈ ਦਿੰਦਾ ਹੈ ਤੇ ਉਸ ਕੋਲ ਇਕ \VUgras{ਪੇਂਡੂ ਤੀਵੀਂ} \textsuperscript{(70)}। ਫਿਰ ਇਕ
\VUgras{ਬੁੱਢੀ} \textsuperscript{(71)}, ਜੋ ਆਪਣੀ \VUgras{ਸੋਟੀ} \textsuperscript{(72)}
ਦੇ ਸਹਾਰੇ ਹੈ। ਅੰਤ ਵਿਚ \VUgras{ਚੱਕੀ ਵਾਲਾ} \textsuperscript{(73)} ਆਪਣੀ ਚੌੜੀ
\VUgras{ਨਮਦੇ ਦੀ ਹੈਟ} \textsuperscript{(74)} ਵਿਚ, ਤੇ \VUgras{ਖੜਾਵਾਂ}
\textsuperscript{(75)} ਪਾਈ ਇਕ ਜਵਾਨ ਪੇਂਡੂ ਤੀਵੀਂ, ਜੋ ਆਪਣੇ ਬਾਲ ਨੂੰ ਤੁਰਨਾ ਸਿਖਾ ਰਹੀ
ਹੈ।

%%K t03-c1-24-1 p
24. --- ਦ੍ਰਿਸ਼ ਬਹੁਤ ਜੀਵੰਤ ਹੈ।
\VUsaut{4.0mm}

%%K t03-tit-8 sub
\VUcentre{12.0pt}{0}{\textit{ਦੂਜਾ ਦ੍ਰਿਸ਼।}}
\VUsaut{3.0mm}

%%K t03-tit-9 sub
\VUpk{10.2pt}{40}{ਸਰਬਜਨਕ ਮੇਲਾ।}
\VUsaut{2.4mm}

%%K t03-tit-10 sub
\VUpk{10.2pt}{40}{ਪਾਣੀ ਦੀਆਂ ਖੇਡਾਂ ਤੇ ਬੇੜੀਆਂ ਦੀ ਦੌੜ।}
\VUsaut{3.0mm}

%%K t03-c2-01-1 p
1. --- \VUgras{ਗਰਮੀ} ਆ ਗਈ ਹੈ। \VUgras{ਜੂਨ} (ਜੁਲਾਈ ਜਾਂ \VUgras{ਅਗਸਤ}) ਦੇ ਮਹੀਨੇ
ਦਾ ਤਪਦਾ ਸੂਰਜ ਜਵਾਨਾਂ ਨੂੰ ਨ੍ਹਾਉਣ ਤੇ ਬੇੜੀ ਚਲਾਉਣ ਲਈ ਲਲਚਾਉਂਦਾ ਹੈ।

%%K t03-c2-02-1 p
2. --- ਦਰਿਆ ਦੇ ਸੱਜੇ ਕੰਢੇ ਉੱਤੇ ਮਲਾਹ ਪਾਣੀ ਦੀਆਂ ਖੇਡਾਂ ਤੇ ਬੇੜੀਆਂ ਦੀ ਦੌੜ ਵਿਚ ਉਤਰਨ
ਲਈ ਕੱਪੜੇ ਲਾਹ ਰਹੇ ਹਨ।

%%K t03-c2-03-1 p
3. --- ਉਨ੍ਹਾਂ ਵਿਚੋਂ ਇਕ ਆਪਣੇ \VUgras{ਜੁੱਤੇ} \textsuperscript{(1)} ਲਾਹ ਰਿਹਾ ਹੈ; ਉਹ
ਉਨ੍ਹਾਂ ਦੇ ਤਸਮੇ (ਬਟਨ) ਖੋਲ੍ਹ ਰਿਹਾ ਹੈ। ਉਸ ਦੇ ਦੋ ਸਾਥੀ ਪਹਿਲਾਂ ਹੀ ਤਿਆਰ ਹਨ। ਹੁਣ
ਉਨ੍ਹਾਂ ਉੱਤੇ ਸਿਰਫ਼ ਆਪਣੀ \VUgras{ਬਨੈਣ} \textsuperscript{(2)} ਤੇ
\VUgras{ਤੈਰਨ ਦਾ ਜਾਂਘੀਆ} \textsuperscript{(3)} ਹੈ। ਉਹ ਆਪਣੇ ਸਾਥੀਆਂ ਦੀ ਉਡੀਕ ਕਰਦੇ
ਗੱਲਾਂ ਕਰ ਰਹੇ ਹਨ। ਉਹ ਦੂਜੇ ਕੰਢੇ ਉੱਤੇ ਲੱਗੇ ਮੇਲੇ ਤੇ ਜਸ਼ਨ ਦੀ ਗੱਲ ਕਰ ਰਹੇ ਹਨ।

%%K t03-c2-04-1 p
4. --- ਉਨ੍ਹਾਂ ਕੋਲ ਭੁੰਜੇ ਉਨ੍ਹਾਂ ਦੇ ਸਾਥੀਆਂ ਦੇ \VUgras{ਕੱਪੜੇ} ਪਏ ਹਨ: ਇਕ
\VUgras{ਜੋੜੀ} \VUgras{ਬੂਟ} \textsuperscript{(4)}, \VUgras{ਜੁੱਤੇ}
\textsuperscript{(5)}, \VUgras{ਜੁਰਾਬਾਂ} \textsuperscript{(6)}, \VUgras{ਨਮਦੇ}
\textsuperscript{(7)} ਤੇ \VUgras{ਤੀਲਿਆਂ} \textsuperscript{(8)} ਦੀਆਂ ਹੈਟਾਂ ਤੇ ਇਕ
\VUgras{ਟਾਈ} \textsuperscript{(9)}।

%%K t03-c2-05-1 p
5. --- ਇਕ ਜਵਾਨ ਨੇ ਆਪਣਾ \VUgras{ਕੋਟ} \textsuperscript{(10)} ਸਿਆਣਪ ਨਾਲ ਤਹਿ ਕੀਤਾ ਹੈ।
ਉਹ ਉਸ ਨੂੰ ਇਕ ਤੌਲੀਏ ਉੱਤੇ ਰੱਖ ਰਿਹਾ ਹੈ, ਜਿਸ ਵਿਚ ਉਸ ਦਾ ਗੁਆਂਢੀ ਥੋੜੀ ਦੇਰ ਪਿੱਛੋਂ
ਉਨ੍ਹਾਂ ਦੋਹਾਂ ਦੇ \VUgras{ਸੂਟ} ਬੰਨ੍ਹ ਦੇਵੇਗਾ।

%%K t03-c2-06-1 p
6. --- ਛੇਵਾਂ ਮੁੰਡਾ ਪਿੱਛੇ ਰਹਿ ਗਿਆ ਹੈ। ਉਸ ਦੇ ਸਿਰ ਉੱਤੇ ਹਾਲੇ ਵੀ ਉਸ ਦੀ
\VUgras{ਟੋਪੀ} \textsuperscript{(11)} ਹੈ। ਉਸ ਨੇ ਨਾ ਆਪਣੀ \VUgras{ਕਮੀਜ਼}
\textsuperscript{(12)} ਲਾਹੀ ਹੈ, ਨਾ ਆਪਣਾ \VUgras{ਕਾਲਰ} \textsuperscript{(13)}, ਨਾ
ਆਪਣੇ \VUgras{ਕਫ਼} \textsuperscript{(14)}। ਉਹ ਹੱਥ ਵਿਚ ਆਪਣੀ \VUgras{ਵਾਸਕਟ}
\textsuperscript{(17)} ਚੁੱਕੀ ਹੋਈ ਹੈ ਤੇ ਹਾਲੇ ਵੀ ਆਪਣੀ \VUgras{ਪਤਲੂਨ}
\textsuperscript{(16)} ਤੇ ਆਪਣੀਆਂ \VUgras{ਗੈਲਸਾਂ} \textsuperscript{(15)} ਪਾਈ ਹੋਈ ਹੈ।
ਉਹ ਅਗਲੀ ਦੌੜ ਲਈ ਜ਼ਰੂਰ ਤਿਆਰ ਨਾ ਹੋ ਸਕੇਗਾ।

%%K t03-c2-07-1 p
7. --- ਜਵਾਨਾਂ ਕੋਲ \VUgras{ਊਠ-ਕਟਾਰਿਆਂ} \textsuperscript{(18)} ਨਾਲ ਘਿਰੀ ਇਕ ਪਗਡੰਡੀ
ਦਰਿਆ ਦੇ ਕੰਢੇ ਉਤਰਦੀ ਹੈ, ਜਿੱਥੇ ਬੁੱਢਾ ਜੇਮਜ਼ \VUgras{ਕਾਨਿਆਂ} \textsuperscript{(19)}
ਵਿਚਕਾਰ ਸ਼ਾਮ ਨੂੰ ਛੱਡੀ ਜਾਣ ਵਾਲੀ \VUgras{ਆਤਿਸ਼ਬਾਜ਼ੀ} \textsuperscript{(20)} ਤਿਆਰ ਕਰ
ਰਿਹਾ ਹੈ।

%%K t03-tit-11 sub
\VUsaut{4.4mm}
\VUpk{10.2pt}{40}{ਦੌੜ।}
\VUsaut{3.4mm}

%%K t03-c2-08-1 p
8. --- ਦਰਿਆ ਉੱਤੇ ਆਪਣੀਆਂ ਛਤਰੀਆਂ ਦੀ ਛਾਂ ਵਿਚ ਕੁਝ ਬੀਬੀਆਂ ਇਕ \VUgras{ਡੂੰਗੀ}
\textsuperscript{(21)} ਵਿਚ ਨਿਕਲੀਆਂ ਹਨ। ਉਹ ਦੌੜ ਵੇਖ ਰਹੀਆਂ ਹਨ, ਜੋ ਬਹੁਤ ਰੋਮਾਂਚਕ ਹੈ।
ਹਰ \VUgras{ਬੇੜੀ} \textsuperscript{(22)} ਵਿਚ ਚਾਰ \VUgras{ਚੱਪੂ ਵਾਲੇ}
\textsuperscript{(24)} ਪੂਰੀ ਤਾਕਤ ਨਾਲ ਖਿੱਚ ਰਹੇ ਹਨ, ਜਦ ਕਿ \VUgras{ਮਲਾਹ}
\textsuperscript{(26)} \VUgras{ਪਤਵਾਰ} \textsuperscript{(25)} ਫੜੀ ਨਾਜ਼ੁਕ ਬੇੜੀ ਨੂੰ
ਸੇਧ ਦਿੰਦਾ ਹੈ। \VUgras{ਚੱਪੂ} \textsuperscript{(23)} ਤਾਲ ਨਾਲ ਪਾਣੀ ਉੱਤੇ ਪੈਂਦੇ ਹਨ, ਤੇ
ਹੌਲੀਆਂ ਬੇੜੀਆਂ ਚੱਕਰ ਲਿਆ ਦੇਣ ਵਾਲੀ ਚਾਲ ਨਾਲ ਦੌੜਦੀਆਂ ਹਨ।

%%K t03-c2-09-1 p
9. --- ਕੁਝ ਜਵਾਨ ਪੁਲ ਦੇ ਥੰਮ੍ਹ ਕੋਲ \VUgras{ਚਿਕਣੇ ਬਾਂਸ} \textsuperscript{(28)} ਦਾ
ਇਨਾਮ ਜਿੱਤਣ ਵਿਚ ਲੱਗੇ ਹਨ। ਉਨ੍ਹਾਂ ਵਿਚੋਂ ਇਕ ਲਚਕਦੇ ਤਿਲਕਣੇ ਬਾਂਸ ਉੱਤੇ ਸਿਆਣਪ ਨਾਲ
ਸਰਕਦਾ ਹੋਇਆ ਸਿਰੇ ਉੱਤੇ ਬੰਨ੍ਹੀ ਨਿੱਕੀ \VUgras{ਝੰਡੀ} \textsuperscript{(29)} ਤਕ ਪਹੁੰਚਣਾ
ਚਾਹੁੰਦਾ ਹੈ। ਉਸ ਦਾ ਨਿਘਰਿਆ \VUgras{ਵਿਰੋਧੀ} \textsuperscript{(30)} ਪਾਣੀ ਵਿਚ ਡਿੱਗ
ਪਿਆ ਹੈ। ਉਹ ਤੈਰ ਕੇ ਕੰਢੇ ਮੁੜ ਰਿਹਾ ਹੈ।

%%K t03-c2-10-1 p
10. --- ਵੱਡੇ ਮੁੰਡੇ \VUgras{ਪਾਣੀ ਦੀ ਲੜਾਈ} \textsuperscript{(32)} ਖੇਡ ਰਹੇ ਹਨ,
\VUgras{ਘਾਟ} \textsuperscript{(33)} ਕੋਲ। ਇਕ ਵੱਡੀ \VUgras{ਭੀੜ} \textsuperscript{(31)}
\VUgras{ਪੁਲ} \textsuperscript{(27)} ਦੀ \VUgras{ਮੁੰਡੇਰ} ਉੱਤੇ ਝੁਕੀ ਤੇ
\VUgras{ਕੰਢੇ} ਉੱਤੇ ਠਸਾਠਸ ਭਰੀ ਇਹ ਸਾਰੀਆਂ ਖੇਡਾਂ ਵੇਖ ਰਹੀ ਹੈ। ਪਰ ਕੁਝ ਤਮਾਸ਼ਬੀਨ ਮੁੜ ਕੇ
\VUgras{ਚੜ੍ਹਨ ਦੇ ਬਾਂਸ} \textsuperscript{(45)} ਦੀ ਖੇਡ ਵੇਖਦੇ ਹਨ, ਜਾਂ
\VUgras{ਇਨਾਮ-ਵੰਡਾਈ} ਵੱਲ ਜਾਂਦੇ \VUgras{ਸ਼ਹਿਰ-ਮੁਖੀ ਦੇ ਜਲੂਸ} ਨੂੰ ਤੱਕਦੇ ਹਨ।

%%K t03-c2-11-1 p
11. --- ਸਭ ਤੋਂ ਪਹਿਲਾਂ ਇਹ ਰਹੇ ਕੁਝ \VUgras{ਛੋਕਰੇ} \textsuperscript{(35)}, ਜੋ
\VUgras{ਸਾਜ਼ ਵਾਲਿਆਂ} \textsuperscript{(36)} ਦੇ ਅੱਗੇ ਅੱਗੇ ਦੌੜ ਰਹੇ ਹਨ; ਫਿਰ ਆਉਂਦੇ
ਹਨ \VUgras{ਸ਼ਹਿਰ-ਮੁਖੀ} \textsuperscript{(37)}, ਉਨ੍ਹਾਂ ਦੇ ਨਾਇਬ ਤੇ ਨਿੱਕੇ ਸ਼ਹਿਰ ਦੀ
\VUgras{ਸ਼ਹਿਰੀ ਸਭਾ}। ਸਭ ਤੋਂ ਪਿੱਛੇ \VUgras{ਅੱਗ ਬੁਝਾਊ ਵਾਲੇ} \textsuperscript{(38)}
ਤੁਰਦੇ ਹਨ।

%%K t03-tit-12 sub
\VUsaut{5.0mm}
\VUpk{10.2pt}{40}{ਮੇਲਾ।}
\VUsaut{3.6mm}

%%K t03-c2-12-1 p
12. --- \VUgras{ਮੇਲੇ ਦੇ ਮੈਦਾਨ} \textsuperscript{(39)} ਉੱਤੇ ਵੱਡੀ ਭੀੜ ਹੈ। ਲੋਕ ਕਈ
ਤਰ੍ਹਾਂ ਦੇ \VUgras{ਤੰਬੂ} ਵੇਖਣ ਆਉਂਦੇ ਹਨ: \VUgras{ਲੱਕੜ ਦੇ ਘੋੜੇ}
\textsuperscript{(40)}, \VUgras{ਸਰਕਸ} \textsuperscript{(41)}, ਜਿੱਥੇ ਖੇਡ ਸ਼ੁਰੂ ਵੀ ਹੋ
ਚੁੱਕੀ ਹੈ; \VUgras{ਸਰਬਜਨਕ ਨਾਚ} \textsuperscript{(42)}, ਜਿੱਥੇ ਸ਼ਹਿਰ ਦੇ ਜਵਾਨ ਮੁੰਡੇ
ਤੇ ਕੁੜੀਆਂ \VUgras{ਨੱਚਣ} ਦਾ ਸੁਆਦ ਲੈ ਰਹੇ ਹਨ।

%%K t03-c2-13-1 p
13. --- ਸਿਰੇ ਉੱਤੇ ਸਾਨੂੰ \VUgras{ਅਖਾੜਾ} \textsuperscript{(44)} ਵਿਖਾਈ ਦਿੰਦਾ ਹੈ,
ਜਿੱਥੇ ਸੂਰਮਾ ਸਪੇਨੀ \VUgras{ਬਲਦ-ਬਾਜ਼} ਗ਼ੁੱਸੇ ਵਿਚ ਆਏ \VUgras{ਬਲਦ}
\textsuperscript{(45)} ਨਾਲ ਲੜਦਾ ਹੈ; ਫਿਰ \VUgras{ਝੂਲੇ ਵਾਲੀ ਰੇਲ} \textsuperscript{(49)}
ਤੇ \VUgras{ਨਿਸ਼ਾਨੇਬਾਜ਼ੀ ਦਾ ਤੰਬੂ} \textsuperscript{(46)}, ਜਿੱਥੇ ਲੋਕ
\VUgras{ਬੰਦੂਕ} ਚਲਾਉਣ ਤੇ \VUgras{ਨਿਸ਼ਾਨੇ} ਉੱਤੇ ਮਾਰਨ ਦਾ ਅਭਿਆਸ ਕਰਦੇ ਹਨ।

%%K t03-c2-14-1 p
14. --- ਕੋਲ ਹੀ ਇਹ ਰਿਹਾ \VUgras{ਮੇਲੇ ਦਾ ਰੰਗ-ਮੰਚ} \textsuperscript{(47)}, ਜਿੱਥੇ
\VUgras{ਹਾਸ-ਨਾਟ} ਤੇ \VUgras{ਨਾਟਕ} ਖੇਡੇ ਜਾਂਦੇ ਹਨ। ਉਸ ਦੇ ਬਿਲਕੁਲ ਕੋਲ
\VUgras{ਦੈਂਤ} \textsuperscript{(48)} ਤੇ \VUgras{ਬੌਣੇ} ਦਾ ਤੰਬੂ ਹੈ। ਪਹਿਲੇ ਦਾ
\VUgras{ਕੱਦ} ਦੋ ਮੀਟਰ ਪੰਦਰਾਂ ਹੈ; ਦੂਜਾ ਮੁਸ਼ਕਲ ਨਾਲ ਇਕ ਬਾਲ ਜਿੰਨਾ ਉੱਚਾ ਹੈ।

%%K t03-c2-15-1 p
15. --- ਅੰਤ ਵਿਚ ਇਹ ਰਹੀ ਵੱਡੀ \VUgras{ਚਿੜੀਆਘਰ ਦੀ ਗੱਡੀ} \textsuperscript{(50)},
ਆਪਣੇ \VUgras{ਜੰਗਲੀ ਜਾਨਵਰਾਂ} ਨਾਲ — ਆਪਣੇ \VUgras{ਸ਼ੇਰ} \textsuperscript{(51)} ਨਾਲ,
ਜਿਸ ਦੇ \VUgras{ਪੰਜੇ} ਬੜੇ ਤਕੜੇ ਹਨ, ਆਪਣੇ \VUgras{ਬੱਬਰ ਸ਼ੇਰ} \textsuperscript{(52)}
ਨਾਲ, ਜਿਸ ਦੀ \VUgras{ਅਯਾਲ} ਸੰਘਣੀ ਹੈ, ਆਪਣੇ ਦੋ \VUgras{ਜ਼ਿਰਾਫ਼ਾਂ}
\textsuperscript{(53)} ਤੇ ਆਪਣੇ \VUgras{ਹਾਥੀ} \textsuperscript{(54)} ਨਾਲ, ਜੋ ਆਪਣੀ
\VUgras{ਸੁੰਡ} ਇੰਨੀ ਸਿਆਣਪ ਨਾਲ ਕੰਮ ਵਿਚ ਲਿਆਉਂਦਾ ਹੈ।

%%K t03-c2-16-1 p
16. --- \VUgras{ਜਾਨਵਰ ਸਿਧਾਉਣ ਵਾਲਾ} \textsuperscript{(55)}, ਜੋ ਇਨ੍ਹਾਂ ਜਾਨਵਰਾਂ ਨੂੰ
ਸਿਖਾਉਂਦਾ ਹੈ, ਤੰਬੂ ਦੇ ਬੂਹੇ ਉੱਤੇ ਖੜਾ ਹੈ।
\VUsaut{6.0mm}
\VUfilet{22mm}
